
\subsection{ProofOptimizer: LLMs for Proof Simplification}

By focusing on proof simplicity as the target metric, we developed
a training and inference-time pipeline for simplifying Lean proofs.

\subsubsection{Training} \label{sec:expit-rl}

First, we trained a general-purpose Lean model by fine-tuning \texttt{Qwen\allowbreak-2.5\allowbreak-7B\allowbreak-Instruct} on a combination of five tasks: natural language problem solving, Lean 4 code completion, auto-formalization (problems and solutions), formal theorem proving, and tactic/proof state prediction.
Next, we describe our two training recipes: expert iteration and online reinforcement learning.
For both recipes, we generated a data mix using a four-stage pipeline involving problem collection, proof
sketching, theorem extraction and filtering, and proof generation.

\textit{ProofOptimizer-ExpIt: Expert Iteration}:
We leveraged a STaR-like \citep{zelikman2022star} iterative training algorithm to improve our model. At a high level, we started with our base model $\pi_0$ and the collection of 145K proofs $P_0$. At each iteration, we attempted to simplify each proof, trained our model on successful proof simplifications, and used the collection of simplified proofs as seed proofs for the next iteration. More precisely, at each iteration $i$, we did the following:
\begin{enumerate}
\itemsep1pt
    \item \textbf{Sample}: For each proof $x \in P_i$, use $\pi_i$ to sample $4$ simplifications $Y_p \triangleq \{y_x^{1}, y_x^{2}, y_x^{3}, y_x^{4}\} \sim \pi_i(x)$.
    \item \textbf{Filter}: Use the Lean compiler to find the shortest correct simplification $y_x \in \{x\} \cup Y_x$. Create a training dataset of proof simplifications $D_i = \{(x, y_x) \mid \text{len}(y_x) \le 0.8 \cdot \text{len}(x), x \in P_i\}$. The length constraint is designed to encourage the model to learn more substantial simplifications rather than trivial ones. For iterations after the first, as $x$ may have been simplified from a more complex proof $x' \in P_0$, we also add $(x', y_x)$ pairs to $D_i$, which are valid and larger proof simplifications. Also, collect simplified proofs $\pi_{i+1} = \{s_x \mid x \in P_i\}$ for the next iteration.
    \item \textbf{Train}: Fine-tune $\pi_i$ on $D_i$ to get $\pi_{i+1}$.
\end{enumerate}


\textit{ProofOptimizer-RL: Online Reinforcement Learning} \label{sec:rl}
In addition to expert iteration as described in the previous section, we trained a proof optimizer model with online reinforcement learning. Using the same dataset as in expert iteration, the reinforcement learning task consists in producing a valid but shorter proof $y$ for a statement given an initial proof $x$. The reward is defined as the relative shortening $R(x, y) = \frac{|y| - |x|}{|x|}$
if $y$ is valid and $|y| \leq |x|$, and $R(x, y) = 0$ otherwise. We employed an asynchronous variant of the GRPO algorithm \citep{shao2024grpo} with advantage $A_i = R_i - \frac{1}{k}\sum_{j \leq k} R_j$
baselined with the average reward of $k = 8$ samples, no advantage normalization by standard deviation \citep{liu2025understandingr1zeroliketrainingcritical}, no KL regularization, and omitting sequences with zero advantage.

\subsubsection{Inference-Time} \label{subsec:inference-time-algorithms}
At inference time, we first implemented a symbolic linter that removes extraneous tactics via Lean's \texttt{linter.\allowbreak{}unusedTactic} linter, which detects tactics that do not change the proof state and provides messages like \texttt{'norm\_num' tactic does nothing}.
Then, we performed iterative proof shortening. For a given proof, we sampled $k$ candidate shortenings and took the shortest correct one. Then, we sampled $k$ shortenings of the new proof, took the shortest correct one -- and so on.
